\documentclass[book.tex]{subfiles}

\begin{document}
% Julia Ling Paper can be covered alongside

% This chapter primarily motivates the use of ML in the context of turbulent flow within vehicle aerodynamics
\chapter{Modeling Turbulent Mixing}

\label{chap:ml_turbulence}
\noindent\rule{11cm}{0.4pt} \\
\textbf{Prerequisites:} \autoref{chap:qft_basics}  \\
\textbf{Difficulty Level:} ** \\
\noindent\rule{11cm}{0.4pt}
\vspace{5mm}

%"Film cooling is a critical technology used to cool gas turbine components, and is particularly important in aerospace applications. Designing those systems requires numerical simulations that predict the resulting temperatures, but for industrial purposes only low fidelity Reynolds-averaged Navier-Stokes (RANS) simulations are feasible. RANS simulations depend on models for the effects of turbulence, and existing models are not terribly accurate in film cooling flows. In this talk, I will discuss the possibility of employing machine learning tools to generate improved models for the turbulent scalar flux that could be used in commercial solvers. The framework consists of running well validated high fidelity simulations that resolve the turbulence in a few simple flows, and then using data driven algorithms to learn the correct behavior of the turbulence in a way that can be applied in low fidelity simulations. A simple model, based on random forest regressors, is discussed first. A more complex model based on deep neural networks is presented subsequently. Both models are designed to be numerically stable, invariant to coordinate frame transformations, and readily applicable on any complex, 3D mesh. Our work shows that these models are promising alternatives for turbulent mixing in the Reynolds-averaged context."


In this chapter we'll cover machine learning methods that are promising for fluid turbulence modeling. Specially, we'll present a Reynolds averaged turbulence modelling approach using deep neural networks with embedded invariance \cite{milani2021turbulent}, \cite{milani2020generalization}, \cite{milani2019enriching}. Such techniques could prove useful for industrial processes such as film cooling, a technology used to cool gas turbine components. In standard practice, only low fidelity numerical simulations can be used to design these systems, but new machine learned approaches could dramatically accelerate these applications. The core idea is to run a few expensive simulations to gain data on the turbulence for simple flows, and then use this data to train a downstream machine learning model that is invariant to coordinate transformations, and applicable to any choice of meshing.


\section{Turbulence Modeling}
Reynolds-averaged Navier–Stokes (RANS) models are widely used because of
their computational tractability. There exists significant demand for improved Reynolds-averaged Navier–Stokes
(RANS) turbulence models that are informed by and can represent a richer set of
turbulence physics. The approach we are going to look at in this chapter uses a deep neural network to
learn a model for the Reynolds stress anisotropy tensor from high-fidelity simulation
data. A novel neural network architecture is proposed which uses a multiplicative
layer with an invariant tensor basis to embed Galilean invariance into the predicted
anisotropy tensor. It is demonstrated that this neural network architecture provides
improved prediction accuracy compared with a generic neural network architecture that
does not embed this invariance property. The Reynolds stress anisotropy predictions
of this invariant neural network are propagated through to the velocity field for two
test cases. For both test cases, significant improvement versus baseline RANS linear
eddy viscosity and nonlinear eddy viscosity models is demonstrated.


%\section{Machine Learning Approach}

%\section{Architecture}

%\section{Results}

%\section{Outlook}


\end{document}